\documentclass[11pt,a4paper]{article}

\usepackage{amsmath,amssymb,amsthm}
\usepackage{hyperref}
\usepackage{microtype}
\usepackage{geometry}
\usepackage{doi}
\geometry{margin=2.5cm}

\newtheorem{theorem}{Theorem}[section]
\newtheorem{proposition}[theorem]{Proposition}
\newtheorem{lemma}[theorem]{Lemma}
\newtheorem{corollary}[theorem]{Corollary}
\theoremstyle{definition}
\newtheorem{definition}[theorem]{Definition}
\newtheorem{hypothesis}[theorem]{Hypothesis}
\theoremstyle{remark}
\newtheorem{remark}[theorem]{Remark}

\newcommand{\Heis}{\mathrm{Heis}_3}
\newcommand{\HeisR}{\mathrm{Heis}_3(\mathbb{R})}
\newcommand{\HeisZ}{\mathrm{Heis}_3(\mathbb{Z}/q\mathbb{Z})}
\newcommand{\HeisZZ}{\mathrm{Heis}_3(\mathbb{Z})}
\newcommand{\dCC}{d_{\mathrm{CC}}}
\newcommand{\DeltaH}{\Delta_H}
\newcommand{\LPi}{L_\Pi}
\newcommand{\Gq}{G_q}
\newcommand{\Sq}{S_q}
\newcommand{\Bn}{B_n}
\newcommand{\Sn}{S_n}
\newcommand{\Cq}{\mathcal{C}_q}
\newcommand{\Eq}{\mathcal{E}_q}
\newcommand{\BI}{\mathrm{BI}}
\newcommand{\Dhom}{D_{\mathrm{hom}}}

\title{BFS Shell Stratification and the Emergence of\\
Four-Dimensional Lorentzian Geometry\\[6pt]
{\large Q5b of the Cosmochrony Spectral Geometry Programme}}

\author{J\'{e}r\^{o}me Beau\\[4pt]
\small Independent Researcher\\
\small \texttt{jerome.beau@cosmochrony.org}}

\date{2026}

\begin{document}

  \maketitle

  \begin{abstract}
    The companion paper Q5a establishes, conditionally, that the admissibility forms $\Eq$
    on the Weil representation of $\HeisZ$ converge in the Mosco sense to a second-order
    operator $\LPi = -A\partial_x^2$ on $L^2(\mathbb{R})$, which is described as the
    ``one-dimensional operator-theoretic shadow'' of the homogeneous four-dimensional
    structure of $\HeisR$.
    The present paper completes the geometric reading deferred by Q5a.
    We establish three results.
    First, we show that the BFS shell stratification of $\Gq = \HeisZ$ converges, in
    the pre-saturation regime, to the Carnot--Car\-a\-th\'{e}o\-do\-ry sphere foliation of
    $\HeisR$; the homogeneous dimension $\Dhom = 4$ (Bass--Guivarc'h) endows the limiting
    geometry with the spectral and volume-growth properties of a four-dimensional space
    (Theorem~\ref{thm:carnot}).
    Second, we identify $\LPi$ as the image, under the Schr\"{o}dinger representation,
    of the kinetic sector of the sub-Riemannian Laplacian $\DeltaH$ on $\HeisR$,
    and extract an effective co-metric tensor from the principal symbol of the full
    effective operator (Theorem~\ref{thm:metric}).
    The lifting hypothesis~[H-lift] on which this result was originally conditional
    has since been proved in Q9~\cite{Beau2026q9}, making Theorem~\ref{thm:metric}
    unconditional.
    The metric coefficients have since been fully determined: $A_H = 2$ by
    Q10~\cite{Beau2026q10}, $A_z = 2$ by Q8~\cite{BeauQ8}, and $A_\tau = 2$ by
    Q11~\cite{Beau2026q11}, giving the isotropic Lorentzian co-metric
    $g^{\mu\nu} = 2\,\eta^{\mu\nu}$ with no remaining free parameter.
    Third, the Born--Infeld admissibility constraint forces this metric to carry Lorentzian
    signature $(-,+,+,+)$ (Theorem~\ref{thm:lorentz}, following~\cite{Beau2026c}).
    Each result is given with an explicit status: structural, unconditional (Q9), or
    imported from~\cite{Beau2026c}.
  \end{abstract}

  \paragraph{Keywords.}
    sub-Riemannian geometry, Heisenberg group, Carnot groups, spectral geometry,
    Mosco convergence, non-injective projection, admissibility, Born–Infeld constraint,
    Lorentzian signature, emergent spacetime, principal symbol, co-metric tensor

    \clearpage
    \tableofcontents

    \clearpage
  \section{Introduction}
  \label{sec:intro}

  \subsection{Context and what Q5a establishes}

    The Foundation paper~\cite{Beau2026m} derives, from four axioms governing admissible
    non-injective transitions, that the symmetry group of the admissible fibre $F_n$ is
    isomorphic to $\HeisZ$ for some prime $q$, and that $F_n$ carries the Weil representation
    $V_\rho$ on $L^2(\mathbb{Z}/q\mathbb{Z})$.
    The Cosmochrony white paper~\cite{Cosmochrony} identifies three qualitative convergences
    expected to hold in the large-$q$ limit:
    \begin{enumerate}
      \item[(i)] The group $\Gq = \HeisZ$ and its representation space $\Cq \simeq L^2(\mathbb{Z}/q\mathbb{Z})$
      approach $\HeisR$ and $L^2(\mathbb{R})$.
      \item[(ii)] The BFS shell structure of $\Gq$ approaches a continuous foliation, with
      polynomial ball growth $|\Bn| \sim c\,n^4$ inducing a four-dimensional effective geometry.
      \item[(iii)] The discrete admissibility filter $\Pi_q$ approaches a second-order differential
      constraint, recovering the Born--Infeld Lagrangian in the continuum.
    \end{enumerate}
    The companion paper Q5a~\cite{Beau2026q5a} addresses convergences (i) and (iii).
    It establishes, conditionally on hypotheses [H1], [H2], [H-w], [H-E1] and Conjecture~C,
    three results: a Hilbert inductive limit $\varinjlim_q \Cq \simeq L^2(\mathbb{R})$ (T1);
    convergence of rescaled Weil generators to metaplectic generators (T2); and Mosco
    convergence of the admissibility forms $\Eq$ to a second-order operator
    $\LPi = -A\partial_x^2$ on $L^2(\mathbb{R})$ (T3).
    Q5a states explicitly:
    \begin{quote}
      ``The operator $\LPi$ is the one-dimensional operator-theoretic shadow of the homogeneous
      four-dimensional structure of $\HeisR$ (Bass--Guivarc'h, $D = 4$).
      The identification of $\LPi$ with a four-dimensional effective geometry, the extraction
      of a metric tensor from its symbol, and the selection of the Lorentzian signature by the
      Born--Infeld admissibility condition are the subject of Q5b.''
    \end{quote}
    The present paper delivers this geometric reading.
    Convergence~(ii) and the associated structural results are the primary content.
    The results combine classical geometric convergence theorems, structural consequences
    of the Cosmochrony framework, and conditional identifications stated explicitly as
    hypotheses; each result carries an unambiguous status label throughout.

  \subsection{Scope and main results}

    The paper is organised around three contributions:

    \medskip
    \noindent\textbf{Carnot convergence (Theorem~\ref{thm:carnot}; structural).}
    We show that the BFS stratification $\{\Sn\}$ of $\Gq$ converges, in the pre-saturation
    regime $n \ll q$, to the Carnot--Car\-a\-th\'{e}o\-do\-ry sphere foliation of $\HeisR$.
    The homogeneous dimension $\Dhom = 4$ (Bass--Guivarc'h~\cite{Bass1972,Guivarch1973})
    gives the limiting geometry the spectral dimension of a four-dimensional space.
    This result is structural: it follows from the theory of asymptotic cones of nilpotent
    groups (Pansu~\cite{Pansu1989}) and does not require any new hypothesis.

    \medskip
    \noindent\textbf{Metric extraction (Theorem~\ref{thm:metric}; proved unconditionally
  via Q9~\cite{Beau2026q9}).}
    Under Hypothesis~[H-lift] (Definition~\ref{hyp:lift}, proved in Q9), $\LPi = -A\partial_x^2$ is
    the image, under the Schr\"{o}dinger representation, of the kinetic sector of the
    sub-Riemannian Laplacian $\DeltaH = \tilde X^2 + \tilde Y^2$ on $\HeisR$, restricted
    to the admissible sector.
    The full effective operator on the four-dimensional space $\mathbb{R}_\tau \times \HeisR$
    has a principal symbol $A^{\mu\nu} k_\mu k_\nu$ from which the effective metric tensor
    $g^{\mu\nu} \propto A^{\mu\nu}$ is read off.

    \medskip
    \noindent\textbf{Lorentzian signature (Theorem~\ref{thm:lorentz}; conditional on T5 and
  citing~\cite{Beau2026c}).}
    The Born--Infeld admissibility constraint forces $g^{\mu\nu} = \mathrm{diag}(-1,+1,+1,+1)$
    in homogeneous isotropic admissible regimes.

  \subsection{Notation and conventions}

    Throughout, $q$ denotes a prime and $\Gq = \HeisZ$ the discrete Heisenberg group.
    The BFS ball of radius $n$ is $\Bn = \{g \in \Gq : d(e,g) \leq n\}$ and the BFS shell
    is $\Sn = \Bn \setminus B_{n-1}$.
    The sub-Riemannian Laplacian on $\HeisR$ is denoted $\DeltaH$.
    The limit operator from Q5a is $\LPi = -A\partial_x^2$ acting on $L^2(\mathbb{R})$,
    where $A > 0$ is the admissibility weight from Q5a Hypothesis~[H-w].
    The principal symbol of a differential operator $P$ of order $m$ is
    $\sigma_m(P)(x,\xi)$.
    Status labels follow the convention of~\cite{Beau2026q5a}: ``proved'' denotes a
    complete proof in the cited companion; ``structural'' denotes a consequence of classical
    results cited explicitly; ``conditional'' denotes a result following from a named
    hypothesis.

  \section{Sub-Riemannian Structure of $\HeisR$}
  \label{sec:heis}

  \subsection{The Heisenberg group as a Carnot group}

    The real Heisenberg group $\HeisR$ is the Lie group $\mathbb{R}^3$ equipped with the
    group law
    \begin{equation}
      (x,y,z) \cdot (x',y',z') = \bigl(x+x',\; y+y',\; z + z' + \tfrac{1}{2}(xy' - x'y)\bigr).
      \label{eq:heis-law}
    \end{equation}
    Its Lie algebra $\mathfrak{h}$ is spanned by left-invariant vector fields
    \begin{equation}
      \tilde X = \partial_x + \tfrac{y}{2}\,\partial_z, \qquad
      \tilde Y = \partial_y - \tfrac{x}{2}\,\partial_z, \qquad
      \tilde Z = \partial_z,
      \label{eq:generators}
    \end{equation}
    satisfying $[\tilde X, \tilde Y] = -\tilde Z$ and $[\tilde X, \tilde Z] =
  [\tilde Y, \tilde Z] = 0$.
    The Lie algebra admits a two-step stratification:
    \begin{equation}
      \mathfrak{h} = \mathfrak{g}_1 \oplus \mathfrak{g}_2, \qquad
      \mathfrak{g}_1 = \mathrm{span}\{\tilde X, \tilde Y\}, \qquad \mathfrak{g}_2 =
      \mathrm{span}\{\tilde Z\}.
    \end{equation}
    The horizontal distribution $\mathcal{H} \subset T\HeisR$ is the left-invariant rank-2
    distribution generated by $\mathfrak{g}_1$; it satisfies the H\"{o}rmander bracket
    condition $\mathfrak{g}_1 + [\mathfrak{g}_1, \mathfrak{g}_1] = \mathfrak{h}$,
    ensuring that any two points are connected by a horizontal curve (Chow--Rashevskii
    theorem).
    The central direction $\tilde Z = \partial_z$, generated by the commutator $[\tilde X,
  \tilde Y]$, is not horizontal; it spans the second layer $\mathfrak{g}_2$ and carries
    weight~$2$ under the anisotropic dilations~\eqref{eq:dilations}.
    This weight-$2$ status is what makes the central direction sub-dominant in the horizontal
    geometry while contributing a full unit to the homogeneous dimension.

  \subsection{Anisotropic dilations and homogeneous dimension}

    The group $\HeisR$ admits a one-parameter family of anisotropic dilations
    \begin{equation}
      \delta_r(x, y, z) = (r x,\; r y,\; r^2 z), \qquad r > 0,
      \label{eq:dilations}
    \end{equation}
    which are group homomorphisms and assign weights $1$ to $\tilde X$ and $\tilde Y$,
    and weight $2$ to $\tilde Z$.
    The \emph{homogeneous dimension} is defined as
    \begin{equation}
      \Dhom = \sum_i \mathrm{weight}_i \cdot \dim(\mathfrak{g}_i)
      = 1 \cdot 2 + 2 \cdot 1 = 4.
      \label{eq:dhom}
    \end{equation}

    \begin{theorem}[Bass~\cite{Bass1972}, Guivarc'h~\cite{Guivarch1973}]
      \label{thm:bass}
      For the discrete Heisenberg group $\HeisZZ$ with the standard symmetric generating set
      $\{X^{\pm 1}, Y^{\pm 1}\}$, the BFS ball satisfies
      \begin{equation}
        |B_n(\HeisZZ)| \sim C \, n^4 \quad \text{as } n \to \infty,
      \end{equation}
      for a constant $C > 0$.
      The same growth law $|\Bn(\Gq)| \sim C\,n^4$ holds for $\Gq = \HeisZ$ in the
      pre-saturation regime $n \ll q$.
    \end{theorem}

    The exponent $4 = \Dhom$ is the Bass--Guivarc'h homogeneous dimension of $\HeisR$.
    A space with ball growth $|B_r| \sim r^D$ has Hausdorff dimension $D$ with respect
    to the associated metric.
    Consequently, $(\HeisR, \dCC)$ has Hausdorff dimension $4$, matching the spectral
    properties of a four-dimensional physical space (see Section~\ref{sec:carnot}).

  \subsection{Carnot--Car\'{a}th\'{e}odory metric}

    Fix a left-invariant sub-Riemannian metric $g_{\mathcal{H}}$ on $\mathcal{H}$ making
    $\{\tilde X, \tilde Y\}$ orthonormal.
    The Carnot--Car\-a\-th\'{e}o\-do\-ry distance between two points $p, q \in \HeisR$ is
    \begin{equation}
      \dCC(p,q) = \inf_\gamma \int_0^1 \|\dot\gamma(t)\|_{\mathcal{H}}\,dt,
    \end{equation}
    where the infimum is over all horizontal curves $\gamma$ from $p$ to $q$
    (curves whose velocity lies in $\mathcal{H}$ at every point).
    This metric is finite by Chow--Rashevskii and is left-invariant under the group action.
    The balls $B_r^{\dCC}(e) = \{g \in \HeisR : \dCC(e,g) \leq r\}$ satisfy
    $\mathrm{Vol}(B_r^{\dCC}) \sim C\,r^4$ as $r \to \infty$, consistently with~\eqref{eq:dhom}.

  \subsection{The sub-Riemannian Laplacian}

    The \emph{Kohn--Laplacian} (sub-Laplacian) is the second-order differential operator
    \begin{equation}
      \DeltaH = \tilde X^2 + \tilde Y^2
      = \Bigl(\partial_x + \tfrac{y}{2}\partial_z\Bigr)^2 +
      \Bigl(\partial_y - \tfrac{x}{2}\partial_z\Bigr)^2.
      \label{eq:sublaplacian}
    \end{equation}
    It is left-invariant, homogeneous of degree $-2$ under $\delta_r$, and hypoelliptic
    (H\"{o}rmander's theorem).
    Its principal symbol at $p = (x,y,z)$ and covector $\xi = (\xi_x, \xi_y, \xi_z)$ is
    \begin{equation}
      \sigma_2(\DeltaH)(p,\xi)
      = \Bigl(\xi_x + \tfrac{y}{2}\xi_z\Bigr)^2 + \Bigl(\xi_y - \tfrac{x}{2}\xi_z\Bigr)^2.
      \label{eq:symbol-sublaplacian}
    \end{equation}
    This symbol is degenerate in the $\xi_z$ direction (reflecting the sub-Riemannian
    structure) but is positive on horizontal covectors.
    The operator $-\DeltaH$ is essentially self-adjoint on $L^2(\HeisR)$ and its heat kernel
    satisfies
    \begin{equation}
      p_t(e,e) \sim C\,t^{-\Dhom/2} = C\,t^{-2} \quad \text{as } t \to 0^+,
    \end{equation}
    which is the same decay rate as the heat kernel on $\mathbb{R}^4$.
    This is the precise sense in which $(\HeisR, \DeltaH)$ has \emph{spectral dimension~4}.

  \section{BFS Stratification and Carnot Convergence}
  \label{sec:carnot}

  \subsection{Asymptotic cone of the Heisenberg group}

    The following theorem is the cornerstone of the BFS-to-Carnot identification.

    \begin{theorem}[Pansu~\cite{Pansu1989}]
      \label{thm:pansu}
      Let $\Gamma = \HeisZZ$ with word metric $d_w$ associated to the generating set
      $S = \{X^{\pm 1}, Y^{\pm 1}\}$.
      For every $n \geq 1$, define the rescaled metric space
      $(\Gamma, n^{-1}d_w)$.
      As $n \to \infty$, this sequence of metric spaces converges in the Gromov--Hausdorff
      sense to $(\HeisR, \dCC)$.
      More precisely, for any $R > 0$ and $\varepsilon > 0$, there exists $N_0$ such that
      for all $n \geq N_0$, the ball $\{g \in \Gamma : d_w(e,g) \leq nR\}$ equipped with
      the metric $n^{-1}d_w$ is $\varepsilon$-close to $B_R^{\dCC}(e)$ in the
      Gromov--Hausdorff distance.
    \end{theorem}

  \subsection{Application to the finite Heisenberg group}

    The finite group $\Gq = \HeisZ$ differs from $\HeisZZ$ only through the identification
    $k \equiv k + q$ in each coordinate direction.
    In the pre-saturation regime $n \ll q$, no such identification is active: the BFS ball
    $\Bn(\Gq)$ is isomorphic, as a metric ball, to $\Bn(\HeisZZ)$.

    \begin{theorem}[Carnot convergence; structural]
      \label{thm:carnot}
      Fix $R > 0$.
      In the double limit $q \to \infty$, $n = \lfloor R q^\alpha \rfloor$ with $\alpha < 1$,
      the rescaled BFS metric space $(\Gq, n^{-1}d_w)$ restricted to balls of radius $R$
      converges in the Gromov--Hausdorff sense to the Carnot--Car\-a\-th\'{e}o\-do\-ry
      ball $(B_R^{\dCC}(e) \subset \HeisR, \dCC)$.
      The homogeneous dimension $\Dhom = 4$ gives the limiting space the volume growth and
      spectral dimension of a four-dimensional metric space.
    \end{theorem}

    \begin{proof}
      For $n \ll q$, the BFS ball $\Bn(\Gq) \cong \Bn(\HeisZZ)$ (no wraparound in any
      coordinate for $n < q/2$).
      This identification relies on the asymptotic cone description of nilpotent groups
      in the sense of Pansu~\cite{Pansu1989} and on the Bass--Guivarc'h polynomial growth
      formula~\cite{Bass1972,Guivarch1973}.
      Theorem~\ref{thm:pansu} applied to $\Gamma = \HeisZZ$ then gives the Gromov--Hausdorff
      convergence at scale $n$.
      The Hausdorff dimension and spectral dimension of $(\HeisR, \dCC)$ equal $\Dhom = 4$
      by the Bass--Guivarc'h formula~\eqref{eq:dhom} and the heat-kernel asymptotics of
      Section~\ref{sec:heis}.
    \end{proof}

  \subsection{Effective four-dimensional foliation}

    The BFS stratification $\{\Sn\}_{n \geq 0}$ of $\Gq$ is a partition of $\Gq$ by
    ``shells'' at fixed word-metric distance from the identity.
    In the large-$q$ limit, Theorem~\ref{thm:carnot} identifies this stratification with
    the foliation of $(\HeisR, \dCC)$ by CC spheres
    $\partial B_r^{\dCC}(e)$ at increasing radii $r = n/n_{\max}$.

    \begin{remark}[The four-dimensional reading]
      \label{rem:4d}
      In Cosmochrony, the observable spacetime is the product
      $\mathbb{R}_\tau \times \HeisR$, where $\tau$ denotes the projective ordering
      index (BFS depth) from~\cite{Beau2026pto,Cosmochrony}.
      As a smooth manifold, $\mathbb{R}_\tau \times \HeisR \cong \mathbb{R}^4$.
      The product structure is not postulated: it reflects the separation between
      the projective ordering direction~$\tau$ (structural output of the admissibility
      cascade~\cite{Beau2026pto,Cosmochrony}) and the relational spatial structure
      encoded by $\HeisR$ (structural output of the Heisenberg symmetry derivation
      of~\cite{Beau2026m,Beau2026HeisenbergStructure}).
      The BFS shells $\{\Sn\}$ foliate this product: for fixed $\tau = n$, each shell is a
      copy of the CC sphere $\partial B_n^{\dCC}$ in the Heis$_3$ factor.
      The homogeneous dimension $\Dhom = 4$ of $\HeisR$ means that the spatial factor
      $\HeisR$ alone carries the volume growth and heat-kernel decay of a four-dimensional
      space.
      The full effective spacetime is therefore four-dimensional in the topological sense
      (as $\mathbb{R}^4$) and also carries a four-dimensional spectral structure in the
      spatial factor alone.
      This is the structural content of convergence~(ii) from the Cosmochrony white
      paper~\cite{Cosmochrony}.
    \end{remark}

    \begin{remark}[Relation to the O-series data]
The O-series numerical campaign~\cite{Beau2026q5a} confirms $\Dhom = 4$ empirically:
the fitted ball-growth exponent $\hat D$ converges to $4$ monotonically as $q$ increases
across primes $q \in \{11, \ldots, 401\}$ (O9--O25).
The pre-saturation window of width $\Omega(q^{1/2})$ BFS steps (Q5a Window-Depth Theorem)
ensures that the Pansu convergence of Theorem~\ref{thm:carnot} is visible for all
tested primes.
    \end{remark}

  \section{Lifting Hypothesis and Sub-Riemannian Operator}
  \label{sec:lift}

  \subsection{The Schr\"{o}dinger representation and the 1D shadow}

    The sub-Laplacian $\DeltaH$ acts on functions $f(x,y,z)$ of three variables.
    Q5a, however, produces a \emph{one-dimensional} operator $\LPi = -A\partial_x^2$
    acting on $L^2(\mathbb{R})$.
    Bridging these two objects is the central open step of this paper: we identify the
    natural candidate for this bridge (the Schr\"{o}dinger representation), formulate a
    precise hypothesis [H-lift] in Section~\ref{ssec:hlift}, and give structural evidence.
    Hypothesis~[H-lift] has since been proved in Q9~\cite{Beau2026q9}.

    The connection between $\DeltaH$ and one-dimensional operators runs through the
    Schr\"{o}dinger representation.

    For each $\xi \neq 0$, the Schr\"{o}dinger representation
    $\pi_\xi : \HeisR \to \mathcal{U}(L^2(\mathbb{R}))$ is defined by
    \begin{equation}
      (\pi_\xi(x,y,z) f)(k) = e^{2\pi i \xi (z + yk)} f(k + x), \qquad f \in L^2(\mathbb{R}).
      \label{eq:schrodinger}
    \end{equation}
    The infinitesimal generators act as
    \begin{equation}
      d\pi_\xi(\tilde X) = \partial_k, \qquad
      d\pi_\xi(\tilde Y) = 2\pi i \xi k, \qquad
      d\pi_\xi(\tilde Z) = 2\pi i \xi \, \mathrm{Id}.
      \label{eq:generators-rep}
    \end{equation}
    Under this representation, the sub-Laplacian becomes
    \begin{equation}
      d\pi_\xi(\DeltaH) = d\pi_\xi(\tilde X)^2 + d\pi_\xi(\tilde Y)^2
      = \partial_k^2 - (2\pi\xi)^2 k^2.
      \label{eq:harmonic-osc}
    \end{equation}
    This is the quantum harmonic oscillator (up to sign and scaling).
    The two terms correspond to:
    \begin{itemize}
      \item $\partial_k^2$: the \emph{kinetic} sector, acting purely as $-(-\partial_k^2)$;
      \item $-(2\pi\xi)^2 k^2$: the \emph{potential} sector, growing with $k^2$.
    \end{itemize}
    The operator $\LPi = -A\partial_x^2$ produced by Q5a corresponds to (a multiple of)
    the kinetic sector, in the variable $k \equiv x$.

  \subsection{Admissibility filter and sector dominance}

    The admissibility filter $\Pi_q$ of Q5a selects a \emph{low-energy} subspace of
    $\Cq = L^2(\mathbb{Z}/q\mathbb{Z})$.
    In the large-$q$ limit, this subspace consists of functions $f \in L^2(\mathbb{R})$
    whose Fourier support is concentrated near $|\xi| \ll \xi_{\max} = q/2$.
    In this regime:
    \begin{itemize}
      \item The potential term $(2\pi\xi)^2 k^2$ contributes at energy scale
      $(2\pi\xi)^2 \langle k^2 \rangle$, which is suppressed relative to the kinetic term
      $\|\partial_k f\|^2$ for admissible functions whose $k$-spread is bounded by
      $A_n^{\max} = c_{\mathrm{BI}} / \sqrt{\lambda_n}$ (BI admissibility envelope).
      \item In the admissible sector, $\langle k^2 \rangle \lesssim (c_{\mathrm{BI}})^2 / \lambda_{\min}$
      is uniformly bounded, while the kinetic energy $\|\partial_k f\|^2$ grows with
      the number of independent Weil fingerprints.
    \end{itemize}
    The Mosco convergence result (T3 of Q5a) shows that the admissibility forms
    $\Eq(f,f) \to A\int |\partial_x f|^2\,dx$, capturing only the kinetic sector.
    The potential term is sub-leading in the admissible regime and does not contribute
    to the Mosco limit.
    This dominance is consistent with the coercivity of the Mosco limit form, which
    excludes oscillatory contributions from the potential sector: a coercive Dirichlet
    form must be positive-definite on its domain, whereas the potential term
    $-(2\pi\xi)^2 k^2$ is negative-definite and would destroy coercivity if retained.
    This behaviour is consistent with the Mosco convergence of the admissible forms
    (T3 of Q5a): Mosco convergence suppresses oscillatory contributions and selects
    the dominant coercive quadratic part, leaving precisely the kinetic sector as the
    unique coercive limit.

  \subsection{The lifting hypothesis}
    \label{ssec:hlift}

    \begin{hypothesis}
      [[H-lift]]
      \label{hyp:lift}
      The operator $\LPi = -A\partial_x^2$ obtained in Q5a as the Mosco limit of the
      admissibility forms $\Eq$ is the image, under the Schr\"{o}dinger representation
      $\pi_1$ at unit central character $\xi = 1$, of the kinetic sector of the sub-Laplacian
      $\DeltaH$ on $\HeisR$, restricted to the admissible (low-energy) subspace of
      $L^2(\mathbb{R})$.
      Equivalently, the full sub-Laplacian $\DeltaH$ on $\HeisR$ is the pre-image (under
      $\pi_1$) of the harmonic oscillator $\partial_k^2 - k^2$, and $\LPi$ is its kinetic
      component in the admissible sector.
    \end{hypothesis}

    \begin{remark}
      [Justification of [H-lift]]
      \label{rem:hlift-justification}
      Hypothesis~[H-lift] is supported by three structural observations:
      \begin{enumerate}
        \item Q5a (T2) establishes that the rescaled Weil generators $\hat X_q / \sqrt{q}$,
        $\hat Y_q / \sqrt{q}$ converge strongly to the metaplectic generators
        $d\pi_1(\tilde X) = \partial_k$ and $d\pi_1(\tilde Y) = ik$, which are precisely the
        generators entering $\DeltaH$ via~\eqref{eq:generators-rep}.
        \item The Mosco limit (T3 of Q5a) captures only the kinetic sector, consistently with
        the sector-dominance analysis of Section~\ref{sec:lift}.2.
        This is a structural prediction: the admissibility filter must project out the potential
        term to produce a coercive (not oscillatory) Dirichlet form.
        \item The operator identity $d\pi_\xi(\DeltaH) = \partial_k^2 - (2\pi\xi)^2 k^2$
        shows that the kinetic sector of $\DeltaH$ and $\LPi$ share the same functional form
        $-A\partial_k^2$ up to the overall constant $A$ (which absorbs the central character
        $\xi$ and the admissibility weight from [H-w] of Q5a).
      \end{enumerate}
      The main open step is a direct proof that the Mosco limit of $\Eq$ coincides with the
      kinetic sector of $d\pi_1(\DeltaH)$ modulo the potential contribution, which is
      deferred to open problem Q5b-O1.
      This lifting should be understood as a representation-theoretic correspondence
      rather than a pointwise identification of operators: it asserts that
      $\LPi$ and the kinetic sector of $d\pi_1(\DeltaH)$ define the same quadratic
      form on the admissible subspace of $L^2(\mathbb{R})$, not that they are equal
      as unbounded operators on all of $L^2(\mathbb{R})$.
      This correspondence is understood in a weak (spectral or representation-theoretic)
      sense rather than as a pointwise operator identity; in particular, the domains
      may differ outside the admissible subspace.
    \end{remark}

    \begin{remark}[Scope of the harmonic analysis]
      \label{rem:harmonic-analysis}
      A full derivation of [H-lift] would require the Plancherel formula for $\HeisR$
      (decomposing $L^2(\HeisR)$ as a direct integral of Schr\"{o}dinger representations
      $\pi_\xi$, $\xi \neq 0$), together with a careful tracking of how the admissibility
      filter $\Pi_q$ acts on each fiber $\xi$.
      This analysis is not developed here.
      The present paper formulates [H-lift] as a precise structural hypothesis and provides
      the evidence described in Remark~\ref{rem:hlift-justification}; the complete harmonic
      analysis is deferred to future work.
    \end{remark}

  \section{Principal Symbol and Effective Metric Tensor}
  \label{sec:metric}

  \subsection{The full effective operator on $\mathbb{R}_\tau \times \HeisR$}

    Under [H-lift], the admissibility Dirichlet form in the continuum limit corresponds
    to the kinetic sector of $\DeltaH$ acting on the spatial factor $\HeisR$.
    The complete effective operator on the four-dimensional space
    $\mathbb{R}_\tau \times \HeisR$ incorporates both the projective ordering direction
    $\tau$ and the Heisenberg spatial directions:
    \begin{equation}
      L_{\mathrm{eff}} = -A_\tau \partial_\tau^2 + A_H \DeltaH,
      \label{eq:leff}
    \end{equation}
    where $A_\tau > 0$ is the effective temporal coefficient from the projective ordering
    structure~\cite{Beau2026pto} and $A_H > 0$ is the spatial coefficient from [H-lift]
    (related to the constant $A$ in $\LPi = -A\partial_x^2$).
    The sign structure in~\eqref{eq:leff} reflects the hyperbolic character forced by
    Born--Infeld admissibility (Section~\ref{sec:lorentz}).

    \begin{remark}[Relation to the 1D shadow and naming convention]
On the spatial slice $\{\tau = \text{const}\}$, the operator $L_{\mathrm{eff}}$ reduces to
$A_H \DeltaH$.
In the Schr\"{o}dinger representation at $\xi = 1$, the kinetic sector of $A_H \DeltaH$
maps to $A_H \partial_k^2$, which is $-\LPi = A \partial_x^2$ upon relabeling
$k \to x$ and identifying $A_H = A$.
This confirms the reading of $\LPi$ as the ``one-dimensional shadow'' stated in Q5a.
Note on naming: the companion paper Q6b~\cite{Beau2026q6b} uses the notation $L_\Pi$
to denote the full four-dimensional effective operator on $\mathbb{R}_\tau \times \HeisR$,
which is what the present paper calls $L_{\mathrm{eff}}$.
The notation $L_\Pi$ is reserved here for the one-dimensional Mosco limit operator
of Q5a, consistently with~\cite{Beau2026q5a}.
    \end{remark}

  \subsection{Principal symbol and metric extraction}

    Let $p = (\tau, x, y, z) \in \mathbb{R}_\tau \times \HeisR$ and
    $k = (k_\tau, k_x, k_y, k_z)$ be the dual (momentum) covector.
    The principal symbol of $L_{\mathrm{eff}}$ at $p$ and $k$ is
    \begin{equation}
      \sigma_2(L_{\mathrm{eff}})(p, k) = -A_\tau k_\tau^2 + A_H \sigma_2(\DeltaH)(p, k_x, k_y, k_z),
      \label{eq:symbol-leff}
    \end{equation}
    where $\sigma_2(\DeltaH)$ is given by~\eqref{eq:symbol-sublaplacian}.
    In coordinates, this reads
    \begin{equation}
      \sigma_2(L_{\mathrm{eff}})(p,k)
      = -A_\tau k_\tau^2
      + A_H \Bigl[\Bigl(k_x + \tfrac{y}{2}k_z\Bigr)^2 + \Bigl(k_y - \tfrac{x}{2}k_z\Bigr)^2\Bigr].
      \label{eq:symbol-expanded}
    \end{equation}
    Following the operator--metric correspondence of~\cite{Beau2026c} (see
  also~\cite{Beau2026a}), the effective metric tensor $g^{\mu\nu}(p)$ is read off
    from the principal symbol by
    \begin{equation}
      \sigma_2(L_{\mathrm{eff}})(p,k) = g^{\mu\nu}(p)\,k_\mu k_\nu.
      \label{eq:op-metric}
    \end{equation}
    This defines $g^{\mu\nu}$ as a symmetric tensor on $\mathbb{R}_\tau \times \HeisR$.
    To spell out the identification: the sub-Laplacian $\DeltaH = \tilde X^2 + \tilde Y^2$
    is a sum of squares of horizontal vector fields, hence its principal symbol
    \begin{equation}
      \sigma_2(\DeltaH)(p,\xi) = \sum_{i \in \{x,y\}}
      (\tilde v_i(p) \cdot \xi)^2, \qquad \tilde v_x = \tilde X,\; \tilde v_y = \tilde Y,
      \label{eq:symbol-as-sum}
    \end{equation}
    defines a positive semi-definite quadratic form on $T^*\HeisR$ whose matrix
    \[
      A^{jk}_H(p) = \sum_{i \in \{x,y\}} [\tilde v_i(p)]^j [\tilde v_i(p)]^k
    \]
    is the \emph{horizontal co-metric tensor} of the sub-Riemannian structure.
    The full symbol of $L_{\mathrm{eff}}$ adds the temporal direction with a negative
    coefficient $-A_\tau$, consistently with the hyperbolic structure required
    in Section~\ref{sec:lorentz}:
    \begin{equation}
      g^{\mu\nu}(p) = \mathrm{diag}(-A_\tau) \oplus A_H \cdot A^{ij}_H(p).
      \label{eq:full-cometric}
    \end{equation}
    This identification follows the standard correspondence between second-order
    differential operators and quadratic forms on the cotangent bundle~\cite{Beau2026c,Beau2026a}.
    The metric $g^{\mu\nu}$ is not postulated: it is read off from the operator
    $L_{\mathrm{eff}}$ constructed from the admissibility data, under [H-lift].
    The metric tensor obtained in this way is not a fundamental geometric object
    defined at the level of the substrate~$\chi$, but an effective structure arising
    from the admissible projection of spectral data onto the continuum limit;
    it carries no ontological status beyond this projection.
    This extraction is canonical once the effective operator $L_{\mathrm{eff}}$ is fixed:
    the principal symbol uniquely determines $g^{\mu\nu}$ up to the overall scale
    absorbed into the normalisation of~$c$.

    \begin{theorem}[Metric extraction; conditional on {[H-lift]} and Q5a hypotheses]
      \label{thm:metric}
      Under Hypothesis~[H-lift] and the Q5a hypotheses \textrm{[H1]}, \textrm{[H-w]},
      \textrm{[H-E1]}, \textrm{[C]}, the Mosco limit operator $\LPi = -A\partial_x^2$ lifts
      to the effective operator $L_{\mathrm{eff}}$ on $\mathbb{R}_\tau \times \HeisR$
      given by~\eqref{eq:leff}.
      Its principal symbol defines, via~\eqref{eq:op-metric}, a leading-order effective
      co-metric tensor
      \begin{equation}
        g^{\mu\nu}(p) = \mathrm{diag}(-A_\tau,\; A_H,\; A_H,\; 0_{z\text{-sector (sub-principal)}})
        \label{eq:metric-tensor}
      \end{equation}
      on $\mathbb{R}_\tau \times \HeisR$, where the $z$-component carries sub-Riemannian
      corrections from the Heisenberg structure.
      This tensor should be understood as the leading-order effective co-metric associated
      with the principal symbol; it becomes a non-degenerate metric only after incorporation
      of sub-principal corrections, as discussed in Remark~\ref{rem:zdirection} and
      open problem Q5b-O2.
      In the homogeneous isotropic regime $(x, y) \approx 0$ (near the identity of $\HeisR$),
      the off-diagonal terms in~\eqref{eq:symbol-expanded} vanish and the symbol simplifies to
      $\sigma_2(L_{\mathrm{eff}}) \approx -A_\tau k_\tau^2 + A_H (k_x^2 + k_y^2)$,
      giving a three-dimensional pseudo-Riemannian symbol at leading order.
    \end{theorem}

    \begin{remark}[Unconditional status via Q9]
      \label{rem:q9lift}
      Hypothesis~[H-lift] was the sole conditional input of Theorem~\ref{thm:metric}.
      It has since been proved in Q9~\cite{Beau2026q9} via generator suppression,
      making Theorem~\ref{thm:metric} unconditional (subject only to the Q5a hypotheses
      \textrm{[H1]}, \textrm{[H-w]}, \textrm{[H-E1]}, \textrm{[C]}, which are
      themselves structural properties of the admissibility filter).
    \end{remark}

    \begin{remark}[Status of the $z$-direction]
      \label{rem:zdirection}
      The sub-Riemannian symbol~\eqref{eq:symbol-sublaplacian} is degenerate in $k_z$:
      the $z$-direction (central element $\tilde Z$) does not appear in the principal symbol
      of $\DeltaH$ at leading sub-Riemannian order.
      Its contribution enters via sub-principal terms and is responsible for the distinction
      between the topological dimension $3$ of $\HeisR$ and its Hausdorff dimension $\Dhom = 4$.
      The coefficient $A_z$ of the $z$-direction has since been established:
      Q8~\cite{BeauQ8} proves $A_z = C_{\mathfrak{su}(2)} = 2$ via Casimir rigidity,
      resolving Q5b-O2 at the algebraic level.
      The operator-theoretic sub-principal framework (Q5b-O2) provides an independent
      derivation of this value within the hypoelliptic calculus.
      The Lorentzian signature result (Section~\ref{sec:lorentz}) applies to the non-degenerate
      part of the symbol and holds independently of this subtlety.
      Crucially, the effective four-dimensionality of the geometry arises from the homogeneous
      (Hausdorff and spectral) dimension $\Dhom = 4$ of the Carnot structure, not from a
      non-degenerate Riemannian metric at the principal-symbol level; these two notions of
      dimension must be carefully distinguished.
      The central direction $\tilde Z$ does not carry an independent kinetic term at
      the sub-Riemannian level: its contribution to the principal symbol of $\DeltaH$
      is sub-leading (see~\eqref{eq:symbol-sublaplacian}).
      Its effective dynamical role emerges only after projection, where it couples to
      the ordering parameter~$\tau$ and contributes to the Lorentzian structure
      through the sub-principal corrections addressed in open problem Q5b-O2.
    \end{remark}

    \begin{remark}[Nature of the effective four-dimensionality]
      \label{rem:eff4d}
      The effective four-dimensionality obtained in this paper is not that of a
      non-degenerate pseudo-Riemannian manifold at the principal-symbol level.
      It arises from the homogeneous dimension $\Dhom = 4$ of the underlying Carnot
      structure, which governs both volume growth and spectral scaling, as established
      in Section~\ref{sec:heis}.
      The Lorentzian metric extracted from the principal symbol is therefore an effective
      object, valid after projection of the discrete admissibility structure onto the
      admissible continuum configurations; it does not claim to be a smooth metric at
      every point of $\mathbb{R}_\tau \times \HeisR$ without further hypotheses.
      The transition from the degenerate sub-Riemannian symbol to a non-degenerate
      effective metric is expected to occur through sub-principal corrections to the
      operator $L_{\mathrm{eff}}$, which lift the central direction~$\tilde Z$ and
      generate a full-rank quadratic form at the effective level; this is the precise
      mathematical content of open problem Q5b-O2.
    \end{remark}

  \section{Lorentzian Signature from Born--Infeld Admissibility}
  \label{sec:lorentz}

  \subsection{The signature selection argument}

    The following theorem is a consequence of the Born--Infeld analysis of~\cite{Beau2026c}.
    It is imported from that paper and adapted to the present geometric setting;
    no independent proof is given here.
    We state it for completeness and give the main lines of argument to make the paper
    self-contained.

    \begin{theorem}[Lorentzian signature; conditional on Theorem~\ref{thm:metric},
      citing~{\cite{Beau2026c}}]
      \label{thm:lorentz}
      In homogeneous and quasi-isotropic admissible regimes, the Born--Infeld admissibility
      constraint forces the effective metric defined by~\eqref{eq:op-metric} to have Lorentzian
      signature $(-,+,+,+)$.
    \end{theorem}

    \begin{proof}[Proof sketch, following~{\cite{Beau2026c}} Sections~5--6]
The operator $L_{\mathrm{eff}}$ must define a well-posed initial value problem in the
projective ordering variable $\tau$; this is required by the structural stability of
the admissibility cascade~\cite{Beau2026pto,Cosmochrony}.
The principal symbol $g^{\mu\nu} k_\mu k_\nu$ determines the propagation properties
of effective modes in the continuum limit, so $g^{\mu\nu}$ must be compatible with
hyperbolicity in the sense of H\"{o}rmander.

Four signature classes must be excluded:
\begin{enumerate}
  \item \textbf{Riemannian $(+,+,+,+)$}: produces a positive-definite symbol, hence an
  elliptic operator.
  An elliptic operator in $(\tau, x, y, z)$ implies instantaneous propagation of
  perturbations across the relational network, contradicting the finite Born--Infeld
  capacity bound $A_n^{\max} = c_{\mathrm{BI}} / \sqrt{\lambda_n}$~\cite{Beau2026m,Beau2026c}.
  \item \textbf{Ultra-hyperbolic (two or more time-like directions)}: fails to admit a
  globally well-posed initial value problem and is excluded by spectral stability of the
  effective operator~\cite{Beau2026c}.
  \item \textbf{Degenerate (rank $< 4$)}: excluded at the level of the effective
  metric after incorporation of sub-principal corrections.
  The leading sub-Riemannian principal symbol is degenerate in the central direction
  $k_z$ (Remark~\ref{rem:zdirection}), but the full effective operator, including
  sub-principal contributions from the Heisenberg commutator structure, is expected
  to yield a non-degenerate metric on $\mathbb{R}_\tau \times \HeisR$
  (open problem Q5b-O2).
  The argument here applies to the non-degenerate block $(-A_\tau, A_H, A_H)$
  from the temporal and horizontal directions; the $z$-direction contributes
  positive-definiteness once the sub-principal corrections are incorporated,
  and does not affect the signature count.
  \item \textbf{Parabolic}: implies the existence of a characteristic direction in which
  the Cauchy problem is ill-posed, contradicting the monotonic projective ordering
  of~\cite{Beau2026pto}.
\end{enumerate}
The unique remaining signature compatible with (i)~hyperbolicity of
$\sigma_2(L_{\mathrm{eff}})$, (ii)~a finite characteristic cone, (iii)~monotonic
projective ordering along $\tau$, and (iv)~the bounded Born--Infeld admissibility
constraint, is the Lorentzian signature $(-,+,+,+)$.
The detailed proof is established in~\cite{Beau2026c} Sections~5--6 and depends only
on the bounded-flux constraint and spectral stability, both of which are present in the
Foundation axioms~\cite{Beau2026m}.
After normalising the maximal admissible propagation speed to $c = 1$, the effective
metric is dynamically fixed to $g^{\mu\nu} = \mathrm{diag}(-1,+1,+1,+1)$ in the
homogeneous isotropic regime.
This selection mechanism is intrinsic to the admissibility constraints and does
not rely on any external causal or geometric assumption introduced independently
of the framework.
    \end{proof}

    \begin{remark}[Emergence, not postulation]
      \label{rem:emergence}
      The Lorentzian metric is not introduced as a background structure.
      It is the unique pseudo-Riemannian structure compatible with:
      the Heisenberg group symmetry (giving $\Dhom = 4$),
      the projective ordering (giving the temporal direction $\tau$),
      and the Born--Infeld admissibility bound (excluding all other signatures).
      Minkowski geometry emerges as the unique stable effective description of a maximally
      symmetric admissible regime.
      The emergence of the Lorentzian structure is thus not postulated but results from
      the admissibility constraints governing the projection of unresolved configurations
      onto observable degrees of freedom — the same non-injective projection that generates
      the Heisenberg symmetry group~\cite{Beau2026m,Beau2026HeisenbergStructure} and the BFS shell
      structure~\cite{Cosmochrony}.
      This is the structural content of the claim in~\cite{Cosmochrony} that spacetime
      is a ``large-$q$ limit of the discrete observable space $\mathcal{O}$, with the
      four-dimensional structure inherited from the homogeneous dimension of $G_q$.''
    \end{remark}

  \section{Status Table and Ancillary Questions}
  \label{sec:status}

  \subsection{Status table}

    The following list summarises the logical status of every result in this paper:
    \begin{itemize}
      \item \textbf{BFS $\to$ Carnot convergence} (Theorems~\ref{thm:bass}--\ref{thm:carnot}):
      classical, following Bass--Guivarc'h~\cite{Bass1972,Guivarch1973} and
      Pansu~\cite{Pansu1989}.
      \item \textbf{Operator lifting} [H-lift] (Hypothesis~\ref{hyp:lift}): explicit
      hypothesis, supported by structural evidence from Q5a~T2 but not proved here.
      \item \textbf{Metric extraction} (Theorem~\ref{thm:metric}): structural consequence
      of [H-lift] together with the Q5a hypotheses [H1], [H-w], [H-E1], [C].
      \item \textbf{Lorentzian signature} (Theorem~\ref{thm:lorentz}): derived from the
      Born--Infeld admissibility analysis of~\cite{Beau2026c}, conditional on
      Theorem~\ref{thm:metric}.
    \end{itemize}

    \medskip
    The detailed status of each item is recorded in the table below.

    \begin{center}
      \begin{tabular}{p{3.2cm}p{6.0cm}p{4.8cm}}
        \hline
        Label & Content & Status \\
        \hline
        Thm.~\ref{thm:bass} & Ball growth $|\Bn| \sim C n^4$, $\Dhom = 4$ &
        Proved (Bass, Guivarc'h) \\
        $\Dhom = 4$ & Effective 4D spectral geometry (Carnot) &
        Structural (Carnot geometry) \\
        Thm.~\ref{thm:pansu} & GH convergence to Carnot CC metric &
        Proved (Pansu) \\
        Thm.~\ref{thm:carnot} & BFS shells converge to CC spheres, $\Dhom = 4$ &
        Structural (from above) \\
        {[H-lift]} & $\LPi$ is kinetic sector of $d\pi_1(\DeltaH)$ &
        Hypothesis (Open: Q5b-O1) \\
        Thm.~\ref{thm:metric} & Metric tensor $g^{\mu\nu} \propto A^{\mu\nu}$ from symbol &
        Unconditional (Q9~\cite{Beau2026q9}\,+\,Q5a) \\
        Thm.~\ref{thm:lorentz} & Lorentzian signature $(-,+,+,+)$ &
        Conditional on Thm.~\ref{thm:metric}, citing~\cite{Beau2026c} \\
        \hline
        Q5b-O1 & Prove [H-lift] via Plancherel decomposition &
        Resolved: Q9~\cite{Beau2026q9} \\
        Q5b-O2 & Extend metric to full $4\times 4$ with $z$-direction &
        Resolved: Q8~\cite{BeauQ8} ($A_z=2$) \\
        Q5b-O3 & Identify $A_\tau$, $A_H$ from admissibility data &
        Resolved: Q11~\cite{Beau2026q11} ($A_\tau = 2$; see \S\ref{ssec:resolved-o3}) \\
        \hline
      \end{tabular}
    \end{center}

  \subsection{Ancillary questions and their resolution}

    Three questions arising naturally from the results of this paper lay outside its
    scope; they were addressed by companion papers and are recorded here for completeness.

    \textbf{Q5b-O1} (Proof of [H-lift] via Plancherel decomposition).
    Q9~\cite{Beau2026q9} proves hypothesis~[H-lift] via generator suppression:
    the admissibility filter $\Pi_q$ concentrates on the kinetic sector of
    $d\pi_\xi(\DeltaH)$ unconditionally, making the Mosco lift to $L_{\mathrm{eff}}$
    a theorem rather than a hypothesis.
    Theorem~\ref{thm:metric} is therefore unconditional with respect to~[H-lift].
    The original Plancherel approach outlined here remains of independent interest
    as an alternative derivation.

    \textbf{Q5b-O2} (Extension to a full $4 \times 4$ Lorentzian metric).
    Q8~\cite{BeauQ8} establishes $A_z = C_{\mathfrak{su}(2)} = 2$ via Casimir rigidity
    on $\mathrm{Sym}^2(V_\rho)$ (Q7~\cite{Beau2026q7}): the $\mathfrak{su}(2)$ Casimir
    acts as $2\cdot\mathrm{Id}$ on the spin-$1$ module $\mathrm{Sym}^2(V_\rho)$,
    and this value is inherited by the $z$-direction coefficient of $L_{\mathrm{eff}}$.
    Combined with $A_H = 2$ (resolved below), the effective co-metric takes the form
    $g^{\mu\nu} = \mathrm{diag}(-A_\tau,\, 2,\, 2,\, 2)$.
    The operator-theoretic sub-principal framework (Beals--Greiner calculus, tangent
    groupoid) remains of independent interest as a derivation from within the PDE calculus.

    \textbf{Q5b-O3} (Identification of $A_\tau$ and $A_H$ from admissibility data).
    \label{ssec:resolved-o3}
    Q10~\cite{Beau2026q10} establishes $A_H = 2$ via the spectral universality
    hypothesis~[U] proved in U1~\cite{Beau2026u1}: the admissibility weight satisfies
    $a_q(s)\cdot q^2 \to 2$ as $q\to\infty$ (confirmed numerically for
    $q \in \{29,61,101,151,211\}$, W1~\cite{Beau2026w1}).
    Q11~\cite{Beau2026q11} subsequently identifies $A_\tau = 2$ via a cascade
    internalisation lemma: $\partial_\tau$ is the continuum limit of the cascade
    increment operator $T : \sigma_c(n) \mapsto \sigma_c(n+1)$ by
    Theorem~\ref{thm:carnot}, and SU(2)-equivariance of $T$ under~[U] forces,
    by Schur's lemma and the Casimir normalisation of Q8~\cite{BeauQ8},
    \[
      A_\tau = C_{\mathfrak{su}(2)}\big|_{\mathrm{Sym}^2(V_\rho)} = 2.
    \]
    The effective co-metric is therefore uniquely determined with no free parameter:
    \[
      g^{\mu\nu} = 2\,\eta^{\mu\nu}, \qquad \eta = \mathrm{diag}(-1,+1,+1,+1).
    \]
    Spacetime structure is a derived consequence of the non-injective projection,
    not an input.

  \section*{Acknowledgements}

  The author acknowledges the use of large language models as a supportive tool for
  refining language, structure, and internal consistency during the development of
  this manuscript.

  \bibliographystyle{plain}
  \bibliography{cosmochrony-bibliography,external-refs}

\end{document}
